\title{%
  Qualitative and quantitative methods
}
\author{Daniel Bosk\thanks{%
    This material is available under the Creative Commons 
    Attribution-NonCommercial-ShareAlike (CC-BY-NC-SA) 4.0 international 
    license.
    The material was written with some aid from GitHub Copilot.
}}
\institute{%
  KTH EECS
}

\begin{frame}
  \maketitle
\end{frame}

\mode*

\begin{abstract}
  % What's the problem?
% Why is it a problem? Research gap left by other approaches?
% Why is it important? Why care?
% What's the approach? How to solve the problem?
% What's the findings? How was it evaluated, what are the results, limitations, 
% what remains to be done?

% XXX Summary
\emph{Summary:}
In this seminar we will explore the relationship between quantitative and 
qualitative methods.
We're particularly interested in how they complement each other and why 
\emph{qualitative} methods are so important in cybersecurity.

% XXX Motivation and intended learning outcomes
\emph{Intended learning outcomes:}
This assignment focuses on practice to
\begin{itemize}
  \item \LOrelate;
  \item \LOevaluate;
  \item \LOapply;
  \item \LOcomm.
\end{itemize}

% XXX Prerequisites
\emph{Prerequisites:}
We need basic knowledge of security, corresponding to an introductory course in 
the subject.
We also need a high-level overview of the breadth of research methods used in 
the area of security, corresponding to the overview lecture \enquote{The 
Scientific Method} given in the course.

\end{abstract}

\clearpage

\section{Warm-up}

\begin{frame}
\begin{columns}[T,totalwidth=\linewidth]
  \begin{column}{0.5\linewidth}
    \begin{block}{Quantitative}
      \only<article>{The quantitative papers for comparison are:}
      \begin{itemize}
        \setlength{\itemsep}{0pt}
        \item Shadow: Running Tor in a box
        \item Website Fingerprinting with Website Oracles
        \item Evaluating the Real-World Applicability of WF
        \item Dos and Don'ts of Machine Learning in Security
      \end{itemize}
    \end{block}
  \end{column}
  \begin{column}{0.5\linewidth}
    \begin{block}{Qualitative}
      \only<article>{The qualitative papers for comparison are:}
      \begin{itemize}
        \setlength{\itemsep}{0pt}
        \item Why phishing works
        \item Why Johnny can't encrypt
        \item Comparing the Usability of Cryptographic APIs
      \end{itemize}
    \end{block}
  \end{column}
\end{columns}
  \begin{exercise}[Warm-up]
    \only<article>{Open the seminar with the following prompt. }
    Choose one paper and name one insight that should shape today's
    discussion.
  \end{exercise}
\end{frame}

\paragraph{Warm-up facilitation}
\ltnote{%
  This opener is deliberately diagnostic. Asking students to name one insight
  from one paper lowers the threshold for participation, activates retrieval,
  and gives the room portable ideas to revisit later. The chosen paper and
  insight vary, while the seminar's continuing task of judging security claims
  stays invariant.
}
\begin{description}
  \item[Aim] Move quickly from first reactions to portable insights that can
    be reused throughout the seminar.
  \item[Listen for] Concrete lessons rather than broad likes or dislikes. Good
    examples include scoped simulation validity, the role of priors in false
    positive claims, and the gap between secure designs and secure use.
  \item[Push further] Ask which later paper should be read differently because
    of the insight the student just named.
  \item[Common trap] Letting the room split into isolated paper summaries.
  \item[Possible takeaway] Each paper contributes a result, but also a
    methodological caution that the whole room can carry into later
    discussions.
\end{description}

\begin{frame}
  \begin{question}[Methods in relation]
    \only<article>{Use these questions to compare what each method can and
      cannot tell us.}
    \begin{itemize}
      \item What can formal, quantitative, and qualitative methods each tell
        us that the others cannot?
      \item If a system is secure on paper but insecure in use, which method
        reveals the missing piece?
    \end{itemize}
  \end{question}
\end{frame}

\paragraph{Methods in relation}
\ltnote{%
  This frame is designed as contrast. The object of learning is not the papers
  themselves, but the kind of claim each method can justify. The method varies
  while the security problem stays invariant, helping students discern
  complementarity instead of ranking methods. The second prompt also supports
  epistemic reasoning by making students state what each method cannot justify
  on its own.
}
\begin{description}
  \item[Aim] Establish complementarity before the paper-specific discussions
    begin.
  \item[Listen for] A three-part distinction: formal methods tell what should
    hold under assumptions, quantitative methods tell what happens and how
    often, and qualitative methods tell why it happens in context.
  \item[Push further] Ask what claim each method cannot justify on its own.
  \item[Common trap] Treating methods as rivals to rank instead of tools that
    expose different parts of the same security problem.
  \item[Possible takeaway] Strong security claims emerge when different methods
    constrain each other.
\end{description}

\section{Quantitative methods}

\begin{frame}
  \begin{block}{Papers}
    \only<article>{The quantitative papers in this part are:}
    \begin{itemize}
      \item Shadow: Running Tor in a box
      \item Website Fingerprinting with Website Oracles
      \item Online Website Fingerprinting: Evaluating the Real-World
        Applicability of Website Fingerprinting
      \item Dos and Don'ts of Machine Learning in Security
    \end{itemize}
  \end{block}

  \begin{question}[Cross-paper setup]
    \only<article>{Start by comparing the strength of the claims before
      trusting any one result.}
    \begin{itemize}
      \item Which result looks strongest at first glance?
      \item Which assumption would you check first before trusting it?
    \end{itemize}
  \end{question}
\end{frame}

\paragraph{Quantitative setup}
\ltnote{%
  The quantitative arc is structured as contrast followed by generalization.
  Shadow isolates scoped simulator validation; the two WF papers keep the
  attack family relatively invariant while varying priors, realism, and
  deployment context; \enquote{Dos and Don'ts of Machine Learning in Security}
  then generalizes the reusable validity checks. The recurring
  claim-assumption-limit scaffold reduces extraneous load while keeping
  attention on the critical aspects.
}
\begin{description}
  \item[Aim] Get students to compare the strength of claims against the
    assumptions behind them, not just against each other.
  \item[Listen for] Remarks about metrics, base rates, priors, realism, and
    drift.
  \item[Push further] Ask whether the apparently strongest result is also the
    most externally valid.
  \item[Common trap] Trusting the number before checking how it was produced.
  \item[Possible takeaway] A quantitative result is only as strong as the data,
    metric, and assumptions behind it.
\end{description}

%\section{Users get routed}
%
%\begin{frame}
%  \begin{question}[Users get routed]
%    \begin{itemize}
%      \item What was the research question?
%      \item What methods did they use?
%      \item Why is that a good method?
%      \item Can we use another method?
%    \end{itemize}
%  \end{question}
%\end{frame}
%
%\begin{frame}
%  \begin{question}
%    \begin{itemize}
%      \item Why do we need to use empirical methods in this case?
%      \item What effect does data have on validity?
%    \end{itemize}
%  \end{question}
%\end{frame}

\subsection{Shadow: Running Tor in a box}

\begin{frame}
  \begin{question}[Shadow: Running Tor in a box]
    \only<article>{Use these prompts to discuss what the simulator validation
      actually buys us.}
    \begin{itemize}
      \item If Shadow matches real Tor on selected distributions, what claim
        has actually been validated?
      \item What would still need separate validation before you trusted a
        downstream result built on Shadow?
    \end{itemize}
  \end{question}
\end{frame}

\paragraph{Shadow facilitation}
\begin{description}
  \item[Aim] Land on scoped trust rather than blanket trust.
  \item[Listen for] Calibration against selected metrics, repeated runs,
    reproducibility, and the difference between validating a tool and
    validating every use of that tool.
  \item[Push further] Ask what kind of downstream claim would still require
    fresh real-world data.
  \item[Common trap] \enquote{The simulator looks realistic, so any result
    from it is trustworthy.}
  \item[Connection] Compare this to the realism complaints in the later WF
    papers and to external-validity concerns from \enquote{Dos and Don'ts of
    Machine Learning in Security}.
  \item[Possible takeaway] Validation buys justified trust within scope, not a
    blank cheque.
\end{description}

\begin{frame}
  \begin{question}[Shadow and realism]
    \only<article>{These follow-up questions connect Shadow to broader concerns
      about realism and scope.}
    \begin{itemize}
      \item Which lesson from \enquote{Dos and Don'ts of Machine Learning in
        Security} should make us more cautious about simulator-based results?
      \item When is a simulator the right method, and when does realism become
        the research question itself?
    \end{itemize}
  \end{question}
\end{frame}

\paragraph{Shadow follow-up}
\begin{description}
  \item[Aim] Shift from \enquote{Is the simulator good?} to \enquote{Good for
    which purpose?}
  \item[Listen for] Trade-offs between control and realism, omitted network
    conditions, and the possibility that a method is excellent for mechanism
    questions but weaker for deployment questions.
  \item[Push further] Ask whether the simulator supports a causal explanation,
    an engineering prediction, or both.
  \item[Common trap] Mistaking experimental convenience for external validity.
  \item[Possible takeaway] Better control can weaken external validity if the
    scope is forgotten.
\end{description}

\subsection{WF with Website Oracles}

\begin{frame}
  \begin{question}[WF with Website Oracles]
    \only<article>{Use these prompts to separate a strong modeled result from
      a broad deployment claim.}
    \begin{itemize}
      \item Does this paper show a practical attack, or a strong result under
        strong assumptions?
      \item Which assumptions about priors, popularity, or oracle availability
        carry the result?
    \end{itemize}
  \end{question}
\end{frame}

\paragraph{Website Oracles facilitation}
\begin{description}
  \item[Aim] Separate a strong modeled result from a broad deployment claim.
  \item[Listen for] Priors, popularity distributions, Bayesian updating, and a
    precise account of what the oracle contributes.
  \item[Push further] Ask what has to be true in the world for the low false
    positive claims to matter operationally.
  \item[Common trap] Repeating \enquote{one in a million} style numbers without
    the assumptions that make them possible.
  \item[Connection] Use \enquote{Online Website Fingerprinting} as the obvious
    contrast case.
  \item[Possible takeaway] The paper sharpens the threat model rather than
    settling real-world practicality.
\end{description}

\begin{frame}
  \begin{question}[What is the contribution?]
    \only<article>{Ask these questions to pin down the paper's actual
      contribution.}
    \begin{itemize}
      \item What is the main contribution here: evidence that website oracles
        exist, a new statistical model, or a stronger attack story?
      \item What would \enquote{Online Website Fingerprinting} need to confirm
        before you would call this operationally realistic?
    \end{itemize}
  \end{question}
\end{frame}

\paragraph{Website Oracles follow-up}
\begin{description}
  \item[Aim] Make students name the contribution precisely rather than praising
    it in general terms.
  \item[Listen for] Evidence that oracles exist, the role of side information,
    and the distinction between plausibility and practicality.
  \item[Push further] Ask what additional empirical validation would move the
    claim from plausible to operational.
  \item[Common trap] Calling the attack practical just because the reported
    false positive rate is small.
  \item[Possible takeaway] Attack plausibility and attack practicality are
    different claims.
\end{description}

\subsection{Evaluating WF in Real World}

\begin{frame}
  \begin{question}[Evaluating WF in Real World]
    \only<article>{Use these prompts to surface why real-world evaluation
      changes the result.}
    \begin{itemize}
      \item Why can the same WF idea look strong in the lab and weak online?
      \item Which matters most here: scale, base rates, drift, or noisier
        traffic?
    \end{itemize}
  \end{question}
\end{frame}

\paragraph{Online WF facilitation}
\begin{description}
  \item[Aim] Foreground external validity and drift.
  \item[Listen for] Open-world imbalance, larger monitored sets, concept drift,
    and noisier traffic traces.
  \item[Push further] Ask which earlier WF assumption this paper is really
    attacking.
  \item[Common trap] Thinking that lower numbers necessarily make this a weaker
    paper, rather than a more realistic one.
  \item[Connection] Keep comparing back to the oracle paper so students see the
    same attack family under different assumptions.
  \item[Possible takeaway] Operational context can dominate classifier
    performance.
\end{description}

\begin{frame}
  \begin{question}[Limits of realism]
    \only<article>{These questions invite the same skepticism toward this paper
      that it directs at earlier work.}
    \begin{itemize}
      \item Which earlier WF claim does this paper narrow or correct?
      \item If this paper is more realistic, what important limits still remain
        in its own approach?
    \end{itemize}
  \end{question}
\end{frame}

\paragraph{Online WF follow-up}
\begin{description}
  \item[Aim] Apply the same skepticism to this paper that it applies to other
    WF work.
  \item[Listen for] Its own collection limits, temporal scope, monitored-set
    choices, and remaining simplifications.
  \item[Push further] Ask what further realism would still be missing before a
    deployment claim felt justified.
  \item[Common trap] Treating \enquote{more realistic} as \enquote{fully
    realistic}.
  \item[Possible takeaway] Realism is comparative, not absolute.
\end{description}

\subsection{Dos and Don'ts of Machine Learning in Security}

\begin{frame}
  \begin{question}[Dos and Don'ts of Machine Learning in Security]
    \only<article>{Use these questions to extract reusable checks for any
      ML-for-security paper.}
    \begin{itemize}
      \item When do more data and higher accuracy make a result more
        trustworthy, and when do they only make it more confidently wrong?
      \item If a paper reports only accuracy, what is the first follow-up
        question you would ask before trusting it?
      \item Which pitfall from this paper threatens the other quantitative
        papers most?
    \end{itemize}
  \end{question}
\end{frame}

\paragraph{Dos and Don'ts facilitation}
\begin{description}
  \item[Aim] Extract reusable evaluation questions for any ML-for-security
    paper.
  \item[Listen for] Leakage, sampling bias, spurious correlations, bad metrics,
    drift, and label noise.
  \item[Push further] Ask whether more data would fix the flaw or only amplify
    it.
  \item[Connection] Invite students to name which pitfall threatens Shadow,
    Website Oracles, and Online WF the most.
  \item[Common trap] \enquote{Big dataset plus high accuracy means trustworthy
    science.}
  \item[Possible takeaway] Methodological errors survive scale.
\end{description}


\section{Qualitative methods}

\begin{frame}
  \begin{block}{Papers}
    \only<article>{The qualitative papers in this part are:}
    \begin{itemize}
      \item Why phishing works
      \item Why Johnny can't encrypt
      \item Comparing the Usability of Cryptographic APIs
    \end{itemize}
  \end{block}

  \begin{question}[Cross-paper setup]
    \only<article>{Open the qualitative part by asking which paper best
      explains practical failure.}
    \begin{itemize}
      \item Which paper best explains why a technically sound system can still
        fail?
      \item What kind of evidence did the authors need to reveal that
        mechanism?
    \end{itemize}
  \end{question}
\end{frame}

\paragraph{Qualitative setup}
\ltnote{%
  The qualitative sequence uses the same mechanism-finding question across
  three contexts. The invariant issue is why security fails in practice; what
  varies is where the failure sits: attention and cues, end-user workflow, or
  developer tooling. This helps students see qualitative work as a way of
  explaining breakdown mechanisms in sociotechnical systems rather than as
  merely subjective commentary.
}
\begin{description}
  \item[Aim] Present the qualitative papers as mechanism-finding studies, not
    as soft add-ons.
  \item[Listen for] Mental models, attention, workflow, documentation, social
    context, and design implications.
  \item[Push further] Ask what an outcome metric alone would have missed in
    these cases.
  \item[Common trap] Calling the findings subjective and stopping there.
  \item[Possible takeaway] Qualitative evidence explains why security failures
    happen in practice.
\end{description}

\subsection{Why phishing works}

\begin{frame}
  \begin{question}[Why phishing works]
    \only<article>{Use these prompts to move the discussion away from simple
      user blame.}
    \begin{itemize}
      \item Is phishing success mainly a knowledge problem, an attention
        problem, or a bad decision environment?
      \item What evidence in the paper points away from simple user blame?
      \item Which insight here should change how we design indicators or
        warnings?
    \end{itemize}
  \end{question}
\end{frame}

\paragraph{Phishing facilitation}
\begin{description}
  \item[Aim] Shift the discussion from user blame to the design of the decision
    environment.
  \item[Listen for] Visual heuristics, attention capture, context, workload,
    and misplaced trust in polished interfaces.
  \item[Push further] Ask what concrete interface or workflow change follows if
    the paper is right.
  \item[Common trap] \enquote{Users just need more training.}
  \item[Connection] Compare this paper with \enquote{Why Johnny can't encrypt}
    and the API usability paper as variants of the same sociotechnical failure.
  \item[Possible takeaway] Users act inside an environment that channels what
    they notice and ignore.
\end{description}

\subsection{Why Johnny can't encrypt}

\begin{frame}
  \begin{question}[Why Johnny can't encrypt]
    \only<article>{Use these questions to keep the focus on security in
      use.}
    \begin{itemize}
      \item If the cryptography is strong but ordinary users cannot use it
        correctly, is the system secure in any meaningful sense?
      \item What hidden assumption about users or workflow failed?
      \item Which lesson here also appears in phishing or API usability?
    \end{itemize}
  \end{question}
\end{frame}

\paragraph{Johnny facilitation}
\begin{description}
  \item[Aim] Keep the focus on \enquote{secure in use} rather than only
    \enquote{secure in theory}.
  \item[Listen for] Mismatches between cryptographic soundness and actual user
    workflow, mental-model failures, and interface burden.
  \item[Push further] Ask whether a system that ordinary users cannot operate
    correctly should count as secure.
  \item[Common trap] Privileging technical strength over usable security.
  \item[Connection] Compare the end-user problems here with the developer-side
    problems in the API paper.
  \item[Possible takeaway] Unusable security can nullify strong primitives.
\end{description}

\subsection{Comparing the Usability of Cryptographic APIs}

\begin{frame}
  \begin{question}[Comparing the Usability of Cryptographic APIs]
    \only<article>{These prompts frame developers as users whose tools shape
      security outcomes.}
    \begin{itemize}
      \item If developers finish the task and feel confident, what could still
        be seriously wrong?
      \item What should count as evidence of usable security here: task
        success, confidence, code quality, or something else?
      \item How is this the developer version of the same problem as
        \enquote{Why Johnny can't encrypt}?
    \end{itemize}
  \end{question}
\end{frame}

\paragraph{API usability facilitation}
\begin{description}
  \item[Aim] Make students see developers as users whose tools shape security
    outcomes.
  \item[Listen for] Unsafe defaults, poor documentation, confidence without
    correctness, and the importance of secure examples.
  \item[Push further] Ask which piece of evidence should matter most if the goal
    is secure code rather than pleasant feelings.
  \item[Common trap] Treating task completion or self-reported confidence as
    sufficient evidence of secure API usability.
  \item[Connection] Use this as the developer-facing version of the workflow
    failures in \enquote{Why Johnny can't encrypt}.
  \item[Possible takeaway] API design and documentation are part of the
    security mechanism.
\end{description}

\subsection{Bringing the papers together}

\begin{frame}
  \begin{question}[Bringing the papers together]
    \only<article>{Use these synthesis questions to force cross-method
      transfer.}
    \begin{itemize}
      \item Which insight from a qualitative paper should change how we judge
        one quantitative paper?
      \item Which quantitative insight should change how we would design a
        qualitative follow-up study?
      \item If you could add one method to one paper, what would it be and why?
    \end{itemize}
  \end{question}
\end{frame}

\paragraph{Synthesis facilitation}
\ltnote{%
  The closing discussion is designed for fusion and transfer. Students must
  hold method, claim, evidence, and remaining blind spot together, then carry
  an insight from one paper into another. This interleaving is deliberate: the
  seminar should end with students matching methods to claims, not with
  parallel paper summaries.
}
\begin{description}
  \item[Aim] Force cross-method transfer instead of ending with parallel paper
    summaries.
  \item[Listen for] Explicit pairings across papers, such as a qualitative
    insight changing trust in a quantitative result or a quantitative limit
    motivating a qualitative follow-up.
  \item[Push further] Ask what extra method they would add, and exactly which
    claim it would test.
  \item[Common trap] Ending with \enquote{all methods matter} without saying
    how.
  \item[Possible takeaway] Match the method to the claim, then ask what blind
    spot remains.
\end{description}
